\origpage{138}{ROBERTSON AND SEYMOUR}
\origfig{figures/fig1}{163bp}{\textsc{Figure} 1}

constrained to lie on some bounded number $\rho$ of regions of $G$. We have given a way to solve the Disjoint Connecting Paths problem for this configuration when $\rho = 1$ and when $\rho = 2$. An example due to L. Lov\'asz indicates that the $\rho = 3$ case is less tractable (at least from the point of view of polyhedral combinatorics). When $\rho = 3$, one can sometimes do better with a ``half-integral'' packing of paths than with an integral one. (It follows easily from the results of this paper that such a phenomenon does not occur when $\rho = 1$ or $2$.) In the graph of figure 1, there do not exist three mutually disjoint paths joining $s_i$ and $t_i$ ($i = 1, 2, 3$), respectively; but there do exist six paths, using each vertex at most twice, such that for $i = 1, 2, 3$, two of them join $s_i$ and $t_i$ (Fig. 1).

\sectionhead{Acknowledgment}

Richard Vitray, who has been at different times the student of each of us, has independently worked out a proof of (3.6).

\sectionhead{References}

1. R. M. Karp, On the computational complexity of combinatorial problems, \textit{Networks} \textbf{5} (1975), 45-68. 2. J. F. Lynch, The equivalence of theorem proving and the interconnection problem, \textit{ACM SIGMA Newsletter} \textbf{5:3} (1975), 31-65. 3. N. Robertson and P. D. Seymour, \textit{Graph minors. VII. Disjoint paths on a surface}, submitted for publication.
